\section{What the benchmark teaches us}
\label{sec:discussion}
\paragraph{New patients and new organs are different tests.}
Within-organ, pooled-patient, and unseen-organ evaluation answer different questions. The preliminary
results show a large transfer gap and organ-dependent model rankings, so a pooled score should not be
used as evidence of organ generalization. Training-only panel ranking closes one source of test
information but introduces a real deployment difficulty: targets useful in training organs may be
weakly expressed elsewhere. That difficulty should be measured rather than hidden by test-informed
selection.

\paragraph{Accuracy needs biological context.}
Pearson correlation remains informative when expression varies sufficiently across spots. It becomes
unstable as target variance approaches zero, and excluding undefined correlations alone does not solve
the problem. We therefore advocate a bundle of evidence: per-organ PCC, rank and error metrics,
detection performance, target variance, evaluable-gene counts, and uncertainty. Detection-aware
rescoring is a diagnostic view of model behavior on measurable targets, not a replacement primary
metric and not evidence that low scores are solely metric artifacts.

\paragraph{FiLM helps, but not everywhere.}
Morphology conditioning improves all six representative backbones in the main study in both regimes,
including a strong MLP probe, so its effect is not STFlow-specific. The complete appendix screening
also contains neutral and negative cases. The evidence thus
suggests an architectural hypothesis---repeated morphology injection helps when contextual aggregation
is missing---rather than proving it or supporting the overbroad claim that FiLM always helps. STFlow does not dominate every
baseline or organ, and its local conditioner is not significantly better than the global version.
Cross-organ transfer therefore remains open.

\paragraph{Why the benchmark and intervention belong together.}
A benchmark is most useful when it can do more than rank methods. COAST first exposes a failure under
organ shift, then uses a matched intervention to ask what information helps repair that failure. The
answer is not a universal winner: morphology conditioning improves several different predictors,
while the complete screening shows where it should be redesigned. This combination of stress test and
controlled change is the central contribution promised by the title.

\paragraph{Limitations.}
HEST is imbalanced and contains organs represented by only two slides; STImage slides are strongly
downsampled. Platform and organ are partly confounded in HEST. Frozen UNI features isolate downstream
models but do not measure encoder adaptation. STFlow restricts evaluation to 50 genes. Finally, the
numerical tables in this working draft use optimistic checkpoint selection and are not submission-ready;
all headline results must be rerun with the fixed validation protocol. Our Gene-DML result uses a
feature-matched downstream adaptation rather than the complete original training pipeline, and two
cells currently contain only two seeds. FEAST could not yet be compared reliably: its dense spatial
attention exceeded available GPU memory on large slides and several runs became numerically invalid.
These are unresolved experiments, not negative results; the final paper must include memory-safe,
validation-selected reruns and distinguish faithful reproductions from COAST-compatible adaptations.

\paragraph{A simple reporting standard.}
Future cross-organ studies should disclose patient and organ boundaries; construct targets without
held-out expression; separate known-organ and unseen-organ results; report macro and per-organ scores;
include target detectability and variance; use multiple complementary metrics and uncertainty; and
release fold-level predictions. These practices make evaluation assumptions auditable and comparisons
extensible.

\section{Conclusion}
Histology-to-ST prediction should be evaluated not only by within-cohort correlation, but under explicit
organ shift, leakage-controlled panel construction, and metrics that remain interpretable for sparse
targets. \coast{} provides such a protocol and exposes both a substantial cross-organ generalization
gap and a detectability-dependent limitation of Pearson correlation. A matched intervention across
six representative architectures shows that morphology conditioning can help different predictor
families; complete screening confirms that it is not universal. The benchmark remains far from solved. Releasing the
partitions, panels, predictions, and evaluator can turn cross-organ generalization from an implicit
claim into a reproducible test.
